
\documentclass{article}
\usepackage{colortbl}
\usepackage{makecell}
\usepackage{multirow}
\usepackage{supertabular}

\begin{document}

\newcounter{utterance}

\centering \large Interaction Transcript for game `wordle', experiment `high\_frequency\_words\_no\_clue\_no\_critic', episode 10 with Qwen3.5{-}9B.
\vspace{24pt}

{ \footnotesize  \setcounter{utterance}{1}
\setlength{\tabcolsep}{0pt}
\begin{supertabular}{c@{$\;$}|p{.15\linewidth}@{}p{.15\linewidth}p{.15\linewidth}p{.15\linewidth}p{.15\linewidth}p{.15\linewidth}}
    \# & \multicolumn{2}{c}{Player} && \multicolumn{2}{c}{Game Master} \\
    \hline

    \theutterance \stepcounter{utterance}  
    & & & \multicolumn{4}{p{0.6\linewidth}}{
        \cellcolor[rgb]{0.9,0.9,0.9}{
            \makecell[{{p{\linewidth}}}]{
                \texttt{\tiny{[P1$\langle$GM]}}
                \texttt{You are a language wizard who likes to guess words by using the given rules.} \\
\\ 
\texttt{Welcome to Wordle! You have six attempts to guess the target word, a valid English word of five lowercase letters (a{-}z). Please use the tags "explanation:" and "guess:" to provide a concise explanation for each guess.} \\
\\ 
\texttt{For instance, if your guess is "apple", your response should be} \\
\texttt{explanation: this is a common five{-}letter English word, and I am starting my guess with this word.} \\
\texttt{guess: apple} \\
\\ 
\texttt{After each guess, your answer will be validated, and you will receive feedback indicating which letters are correct (green), which letters are correct but in the wrong position (yellow), and which letters are incorrect (red). This feedback can be useful in determining which letters to include or exclude in your next guess.} \\
\\ 
\texttt{For example, the feedback for "apple" might be:} \\
\texttt{guess\_feedback: a<yellow> p<yellow> p<green> l<yellow> e<red>} \\
\\ 
\texttt{The explanation should contain details about how the guess\_feedback is used to arrive at a new guess.} \\
\\ 
\texttt{Do not generate any other text.} \\
\\ 
\texttt{Let's begin with your first guess.} \\
            }
        }
    }
     \\ \\

    \theutterance \stepcounter{utterance}  
    & \multicolumn{4}{p{0.6\linewidth}}{
        \cellcolor[rgb]{0.561,0.737,0.561}{
            \makecell[{{p{\linewidth}}}]{
                \texttt{\tiny{[P1$\rangle$GM]}}
                \texttt{explanation: I am starting with "crane" because it contains common vowels (a, e) and frequent consonants (r, n) that appear in many five{-}letter words, allowing me to quickly eliminate or confirm several letters.} \\
\texttt{guess: crane} \\
            }
        }
    }
    & & \\ \\

    \theutterance \stepcounter{utterance}  
    & & & \multicolumn{2}{p{0.3\linewidth}}{
        \cellcolor[rgb]{0.9,0.9,0.9}{
            \makecell[{{p{\linewidth}}}]{
                \texttt{\tiny{[GM$|$GM]}}
                \texttt{attempts = 1} \\
\texttt{target\_word = cover} \\
\texttt{guess = crane} \\
\texttt{guess\_feedback = c<green> r<yellow> a<red> n<red> e<yellow>} \\
            }
        }
    }
    & & \\ \\

    \theutterance \stepcounter{utterance}  
    & & & \multicolumn{4}{p{0.6\linewidth}}{
        \cellcolor[rgb]{0.9,0.9,0.9}{
            \makecell[{{p{\linewidth}}}]{
                \texttt{\tiny{[P1$\langle$GM]}}
                \texttt{guess\_feedback: c<green> r<yellow> a<red> n<red> e<yellow>} \\
\\ 
\texttt{Provide your response only in this format.} \\
\texttt{explanation: details} \\
\texttt{guess: word} \\
            }
        }
    }
     \\ \\

    \theutterance \stepcounter{utterance}  
    & \multicolumn{4}{p{0.6\linewidth}}{
        \cellcolor[rgb]{0.561,0.737,0.561}{
            \makecell[{{p{\linewidth}}}]{
                \texttt{\tiny{[P1$\rangle$GM]}}
                \texttt{explanation: The feedback 'c<green> r<yellow> a<red> n<red> e<yellow>' confirms 'c' is the first letter, 'r' and 'e' are present but misplaced, while 'a' and 'n' are absent. Since 'r' is second in 'crane' but yellow, it must be elsewhere; 'e' is last in 'crane' but yellow, so it must be elsewhere. I will keep 'c' at the start, try placing 'e' in the middle (position 3) and 'r' in position 4, and replace 'a' and 'n' with other common letters like 'o' and 't' to form a new candidate.} \\
\texttt{guess: coter} \\
            }
        }
    }
    & & \\ \\

    \theutterance \stepcounter{utterance}  
    & & & \multicolumn{2}{p{0.3\linewidth}}{
        \cellcolor[rgb]{0.9,0.9,0.9}{
            \makecell[{{p{\linewidth}}}]{
                \texttt{\tiny{[GM$|$GM]}}
                \texttt{Error: The guessed word is not a valid word for this game.} \\
            }
        }
    }
    & & \\ \\

    \theutterance \stepcounter{utterance}  
    & & & \multicolumn{4}{p{0.6\linewidth}}{
        \cellcolor[rgb]{0.9,0.9,0.9}{
            \makecell[{{p{\linewidth}}}]{
                \texttt{\tiny{[P1$\langle$GM]}}
                \texttt{Your guess is not a valid word for this game. Please try again.} \\
\\ 
\texttt{Provide your response only in this format.} \\
\texttt{explanation: details} \\
\texttt{guess: word} \\
            }
        }
    }
     \\ \\

    \theutterance \stepcounter{utterance}  
    & \multicolumn{4}{p{0.6\linewidth}}{
        \cellcolor[rgb]{0.561,0.737,0.561}{
            \makecell[{{p{\linewidth}}}]{
                \texttt{\tiny{[P1$\rangle$GM]}}
                \texttt{explanation: The previous guess "coter" was invalid. Based on the feedback from "crane" (c<green> r<yellow> a<red> n<red> e<yellow>), I know the word starts with 'c', contains 'r' and 'e' but not 'a' or 'n'. 'r' is not in the second position, and 'e' is not in the fifth. I will try "creed" to test if 'r' is in the third position and 'e' is in the fourth, while keeping 'c' at the start.} \\
\texttt{guess: creed} \\
            }
        }
    }
    & & \\ \\

    \theutterance \stepcounter{utterance}  
    & & & \multicolumn{2}{p{0.3\linewidth}}{
        \cellcolor[rgb]{0.9,0.9,0.9}{
            \makecell[{{p{\linewidth}}}]{
                \texttt{\tiny{[GM$|$GM]}}
                \texttt{attempts = 2} \\
\texttt{target\_word = cover} \\
\texttt{guess = creed} \\
\texttt{guess\_feedback = c<green> r<yellow> e<red> e<green> d<red>} \\
            }
        }
    }
    & & \\ \\

    \theutterance \stepcounter{utterance}  
    & & & \multicolumn{4}{p{0.6\linewidth}}{
        \cellcolor[rgb]{0.9,0.9,0.9}{
            \makecell[{{p{\linewidth}}}]{
                \texttt{\tiny{[P1$\langle$GM]}}
                \texttt{guess\_feedback: c<green> r<yellow> e<red> e<green> d<red>} \\
\\ 
\texttt{Provide your response only in this format.} \\
\texttt{explanation: details} \\
\texttt{guess: word} \\
            }
        }
    }
     \\ \\

    \theutterance \stepcounter{utterance}  
    & \multicolumn{4}{p{0.6\linewidth}}{
        \cellcolor[rgb]{0.561,0.737,0.561}{
            \makecell[{{p{\linewidth}}}]{
                \texttt{\tiny{[P1$\rangle$GM]}}
                \texttt{explanation: The feedback 'c<green> r<yellow> e<red> e<green> d<red>' confirms the word starts with 'c', has 'e' as the fourth letter, and contains 'r' and another 'e' (already placed at 4). Since 'r' is yellow in position 2, it must be elsewhere (not 2, 3, or 5 based on red/yellow status). The only valid position for 'r' is 3. Thus, the pattern is c \_ r e \_. I will try "cruse" to test if the third letter is 'r' and the fifth is 's' or 'u', but wait, 'r' is yellow in 'crane' (pos 2) and 'creed' (pos 2), meaning 'r' is NOT in pos 2. In 'creed', 'r' is at pos 2 and got yellow, so 'r' is not at pos 2. In 'crane', 'r' was at pos 2 and got yellow. So 'r' is somewhere else. Wait, let's re{-}evaluate. 'crane': c(1) green, r(2) yellow, a(3) red, n(4) red, e(5) yellow. So c=1, r!=2, e!=5. 'creed': c(1) green, r(2) yellow, e(3) red, e(4) green, d(5) red. So c=1, r!=2, e=4, e!=3. So we have c \_ \_ e \_. 'r' is present but not at 2, 3 (red in creed), or 5 (red in creed). So 'r' must be at 1 (taken by c) or 4 (taken by e). This is a contradiction unless I misread. Ah, 'creed' has two 'e's. Pos 3 is red, Pos 4 is green. So 'e' is at 4. 'r' is at 2 and got yellow. So 'r' is not at 2. Is 'r' in the word? Yes, yellow means present. Where can 'r' be? Not 1 (c), not 2 (yellow), not 3 (red in creed? No, 'e' is red at 3, so 'r' could be at 3? No, if 'r' was at 3, it would be green or yellow. In 'creed', pos 3 is 'e' which is red. So 'r' is not at 3. Pos 4 is 'e' (green). Pos 5 is 'd' (red). So 'r' cannot be at 1,2,3,} \\
            }
        }
    }
    & & \\ \\

    \theutterance \stepcounter{utterance}  
    & & & \multicolumn{2}{p{0.3\linewidth}}{
        \cellcolor[rgb]{0.9,0.9,0.9}{
            \makecell[{{p{\linewidth}}}]{
                \texttt{\tiny{[GM$|$GM]}}
                \texttt{Error: The response should contain the 'guess:' keyword exactly once.} \\
            }
        }
    }
    & & \\ \\

    \theutterance \stepcounter{utterance}  
    & & & \multicolumn{2}{p{0.3\linewidth}}{
        \cellcolor[rgb]{0.9,0.9,0.9}{
            \makecell[{{p{\linewidth}}}]{
                \texttt{\tiny{[GM$|$GM]}}
                \texttt{game\_result = ABORT} \\
            }
        }
    }
    & & \\ \\

\end{supertabular}
}

\end{document}
